


Generated by \texttt{python scripts/vf2\_formal\_doc.py}. The models were run again in WSL from a fresh clone of \texttt{round6-v1} with \texttt{scripts/r6\_formal.sh} (ProVerif 2.05 at 300 s, Tamarin 1.12.0 with Maude 3.5.1, one process per model, one thread, 5 GB heap, 300 s). Q1, Q2, ... are the queries of a ProVerif file in the order of the file (for the micro-confirm model: Q1 authenticity, Q2 binding to the flash, Q3 micro-cancel, Q4 reachability, which is false when the event can be reached; for the ratchet models Q1 to Q3 are the three steps). The rerun outputs are in \texttt{formal/rerun\_vf2/}; the stored outputs are in \texttt{formal/out/r6/}. The bounded quarantine exploration (\texttt{python -m crashguard.r6\_quarantine}, FORMAL, bounded: 3 claimed ids, 6 slots, up to 4 attacker frames per slot) was run again as well; its ten keys \texttt{r6.f3b.*} are equal. The \texttt{r6.f.*} keys written from the new outputs have the same values; three of their notes differ only in peak memory (41 against 40 MB).

\textbf{Result: 6 of 6 models give the same verdicts as the stored run.}

{\scriptsize\raggedright
\begin{longtable}{p{0.136\linewidth}p{0.136\linewidth}p{0.136\linewidth}p{0.136\linewidth}p{0.136\linewidth}p{0.136\linewidth}p{0.136\linewidth}}
\toprule
Model & Tool & Stored verdicts & Rerun verdicts & Same & Stored time and memory & Rerun time and memory \\
\midrule
\endhead
R6-1c: micro-confirm authenticity, binding to the flash, micro-cancel & ProVerif 2.05 & Q1=true, Q2=true, Q3=true, Q4=false & Q1=true, Q2=true, Q3=true, Q4=false & yes & 0 s, 10 MB & 0 s, 10 MB \\
R6-1c control: tag without the provisional tag (attack expected) & ProVerif 2.05 & Q1=false, Q2=true, Q3=false & Q1=false, Q2=true, Q3=false & yes & 0 s, 10 MB & 0 s, 10 MB \\
R6-1c: same properties & Tamarin 1.12.0 & executable\_micro=verified, micro\_auth=verified, micro\_binding=verified, cancel\_auth=verified & executable\_micro=verified, micro\_auth=verified, micro\_binding=verified, cancel\_auth=verified & yes & 0 s, 40 MB & 0 s, 42 MB \\
R6-1c control (attack expected) & Tamarin 1.12.0 & executable\_micro=verified, micro\_auth=falsified, micro\_binding=verified, cancel\_auth=falsified & executable\_micro=verified, micro\_auth=falsified, micro\_binding=verified, cancel\_auth=falsified & yes & 1 s, 40 MB & 1 s, 40 MB \\
R6-3a: authenticity of past alerts under a hash ratchet, 3 steps & ProVerif 2.05 & Q1=true, Q2=true, Q3=true, Q4=false & Q1=true, Q2=true, Q3=true, Q4=false & yes & 0 s, 10 MB & 0 s, 9 MB \\
R6-3a: secrecy of past alerts after the current key is revealed, 3 steps & ProVerif 2.05 & Q1=true, Q2=true, Q3=false & Q1=true, Q2=true, Q3=false & yes & 0 s, 9 MB & 0 s, 9 MB \\
\bottomrule
\end{longtable}}

Reading. The micro-confirm properties hold in both tools (acceptance implies the vehicle sent that micro-confirm, bound to a provisional frame of the same vehicle; a micro-cancel is accepted only if the vehicle sent it), the control whose tag omits the provisional tag is attacked in both tools (authenticity and cancel), and the hash ratchet keeps authenticity and secrecy of past alerts within three steps. Quarantine: the rule that counts before verification reaches quarantine of the victim with two frames in slot 0 (334 states visited), the rule that counts after verification visits 7 states and reaches none. The symbolic models cover the message structure stated in \texttt{docs/05\_architecture/FRAME\_LAYOUT.md} section 8 and nothing below it.

